\newcommand{\Oh}{\ensuremath{\mathcal{O}}}
\newcommand{\FIND}[1]{\ensuremath{\mbox{\texttt{find}}(#1)}}
\newcommand{\find}{\texttt{find}}
\newcommand{\union}{\texttt{union}}
\newcommand{\UNION}[1]{\ensuremath{\mbox{\texttt{union}}(#1)}}

\section{Homework 8: Congruence Closure Algorithm and Complexity}

In this assignment, we will demonstrate that the runtime of the congruence closure implemented
using a \emph{union-find} data structure is in $\Oh(n \log n)$. To this end, we will first
introduce the union-find data structure and its operations and then give a precise description
of the congruence closure algorithm, analyzing each step in terms of its complexity.

\subsection{Union-Find Data Structure}

The union-find data structure, also known as the disjoint-set data structure,
efficiently maintains a collection of disjoint sets.
Each set is represented as a tree, with each node pointing to its parent. The
root of each tree is the so-called \emph{representative} of that set.

The union-find data structure supports two main operations:

\begin{itemize}
  \item \FIND{x}: Determines the representative of the set containing the element $x$.
  \item \UNION{x, y}:
Merges the sets with representatives $x$ and $y$ into a single set.
\end{itemize}


\subsubsection{Optimizations}

To improve the efficiency of union-find operations, one can employ two standard heuristics:

\begin{itemize}
  \item \textbf{Path Compression}: During a \find{} operation, we make each node
  on the search path point directly to the root. This flattens the structure of
  the tree, leading to faster subsequent \find{} operations.
  \item
  \textbf{Union by Rank (or Size)}:
  Let the \emph{rank} of a node in the union find structure denote
  an upper bound on the height of the subtree rooted at that node.
  When performing a \union{} operation, we attach the tree
  with the smaller rank (or size) to the root of the tree with the larger rank.
  This prevents the trees from becoming too tall.
\end{itemize}

\subsubsection{Time Complexity}

With these optimizations, both the \find{} and \union{} operations have an amortized
time complexity of $\mathcal{O}(\alpha(n))$, where $\alpha(n)$ is the inverse
Ackermann function. The function $\alpha(n)$ grows extremely slowly and can be
considered a constant (less than 5) for all practical values of $n$.

As this exercise assumes that the union-find data structure is available, we
will not give a detailed proof of the time complexity of the operations.

\subsection{Congruence Closure Algorithm}

The congruence closure algorithm computes the smallest congruence relation over
terms induced by a given set of equations. This is fundamental in automated
reasoning and term rewriting systems.

We first give a formal description of the problem:

\begin{itemize}
  \item \textbf{Signature $\Sigma$}: A set of function symbols $F$,
and predicate symbols $P$, each with an associated arity.
\item \textbf{Terms $T$}: The set of
all terms constructed from $\Sigma$ and a set of variables, using the usual
rules of term formation.
\item \textbf{Equations $E$}: A finite set of equations
${ s_i = t_i \mid 1 \leq i \leq m }$ between terms in $T$.
\item \textbf{Congruence Relation}: An equivalence relation that is compatible with
the function and predicate symbols in $\Sigma$; that is, if $s_i \equiv t_i$ for $1 \leq i
\leq k$, then $f(s_1, \dots, s_k) \equiv f(t_1, \dots, t_k)$ for any $f \in
F$ and $p(s_1, \dots, s_k) \equiv p(t_1, \dots, t_k)$ for any $p \in P$.
\end{itemize}

The goal of the procedure is to compute the smallest congruence relation
$\equiv$ over $\mathcal{T}$ such that all equations in $E$ are satisfied: $s_i
\equiv t_i$ for all $1 \leq i \leq m$.

The algorithm proceeds in three main phases: Initialization, Processing
Input Equations, and Enforcing Congruence Closure. We will describe each phase
and analyze its time complexity.

\subsubsection*{Initialization}
\begin{itemize}
  \item For each term $t \in T$, create a singleton set
${ t }$ in the union-find data structure.
\item Initialize an empty \textbf{signature table} $\mathcal{S}$, typically implemented as a hash table,
to map signatures to representative terms.
\end{itemize}

\paragraph{Time Complexity} $\Oh(n)$, where $n = |T|$.

\subsubsection*{Processing Input Equations}

For each equation $s = t$ in $E$:
\begin{itemize}
\item
Perform $\UNION{\FIND{s}, \FIND{t}}$ to merge the sets
containing $s$ and $t$.
\end{itemize}
\paragraph{Time Complexity} $\Oh(m \cdot \alpha(n))$, where $m = |E|$.

\subsubsection*{Enforcing Congruence Closure}

\begin{itemize} \item Repeat until no new unions are performed:
  \begin{itemize} \item For each compound term $f(t_1, t_2, \dots, t_k) \in T$:
    \begin{itemize}
      \item Let $r_i = \FIND{t_i}$ for $1 \leq i \leq k$.
      \item Construct the \emph{signature} $\sigma = (f, r_1, r_2, \dots, r_k)$.
      \item If $\sigma$ is in $\mathcal{S}$ with associated term $s$:
      \begin{itemize}
        \item Perform $\UNION{\FIND{f(t_1, \dots, t_k)}},
        \FIND{s})$.
      \end{itemize}
      \item Else:
      \begin{itemize}
        \item Add $\sigma$ to $\mathcal{S}$ with associated term $f(t_1, \dots, t_k)$.
      \end{itemize}
    \end{itemize}
  \end{itemize}
\end{itemize}

\paragraph{Time Complexity} We need to consider the needed number of iterations and the time complexity per iterations.

For the time complexity per iteration, finding the representatives for each
compound term is in $\Oh(k \cdot \alpha(n))$, where $k$ is the maximum arity
of functions and predicates in $\Sigma$. As we consider constant-time hash table operations,
constructing the signature and looking it up in $\mathcal{S}$ is in $\Oh(1)$. The union operation
takes $\Oh(\alpha(n))$ time. Thus, the total time per iteration is in $\Oh(k \cdot \alpha(n))$.

For the number of iterations, observe that each \union{} operation reduces the number of distinct
equivalence classes. Since there are at most $n$ terms, there can be at most $n - 1$ unions. However, the
structure of the terms may cause the same unions to be considered multiple times. We claim
that the number of iterations is bounded by $\Oh(\log n)$.

We claim that, with the mentioned optimizations, the maximum rank $R_{\max}$ in a
union-find data structure with $n$ elements is $\lfloor \log_2 n \rfloor$, which
we will show using the following Lemma:

\iffalse
\begin{lemma}
\label{lemma:rank-bound}
In the union-find data structure with union by rank and path compression, the maximum rank of any node is
at most $\lfloor \log_2 n \rfloor$, where $n$ is the number of elements. Additionally, a
tree of rank $r$ contains at least $2^r$ nodes.
\end{lemma}

\begin{proof}
We prove this by an inductive argument, considering what happens if two trees of rank $r$ are united.

\emph{Base case}: A node with rank $0$ is a singleton node, the minimum number of nodes in
a tree of rank 0 is $1 = 2^0$.

\emph{Union of two trees}: Let the rank of a tree be the maximum rank of any
node in that tree.  When two trees of rank $r_1$ and $r_2$ are united and $r_1
\neq r_2$, the resulting tree has rank $\max(r_1, r_2)$ due to the union by rank
heuristic.

Whenever two trees of the same rank $r$ are united, the resulting tree has rank
$r + 1$, and it contains at least $2^{r+1}$ nodes.  Specifically, for a node of
rank $r$, the minimum number of nodes in its subtree is at least $2^r$.
Therefore, the maximum possible rank satisfies:
\[ 2^{R_{\max}} \leq n \implies R_{\max} \leq \log_2 n \]

\end{proof}
\fi

\begin{lemma}
\label{lem:node-bound}
In the union-find data structure with union by rank, the minimum number
of nodes in a tree of rank $r$ is $2^r$.
\end{lemma}

\begin{proof}
  We prove this by an inductive argument, considering that a tree's rank can
  only increase due to \union{} operations.

  \emph{Base case}: For $r = 0$, the tree is a singleton node, and the minimum
  number of nodes is $1 = 2^0$.

  \emph{Inductive step}: Let the rank of a tree be the maximum rank of any
  node in that tree.  When two trees of rank $r_1$ and $r_2$ are united and $r_1
  \neq r_2$, the resulting tree has rank $\max(r_1, r_2)$ due to the union by rank
  heuristic and thus contains at least $2^{\max(r_1, r_2)}$ nodes.

  Whenever two trees of the same rank $r$ are united, the resulting tree has rank
  $r + 1$, as one tree is attached to the other ones root. Thus, the resulting
  tree has at least $2^r + 2^r = 2^{r+1}$ nodes by the induction hypothesis.
\end{proof}

Observe that path compression will only reduce the height of a node's subtree,
and thus has no effect on the presented lower bound.

From Lemma \ref{lem:node-bound}, we immediately obtain the following corollary:

\begin{corollary}
The maximum rank $R_{\max}$ of any node in the union-find data structure with
union by rank is at most $\lfloor \log_2 n \rfloor$.
\end{corollary}

\begin{lemma}
\label{lemma:change-representative}
The maximum number of times the representative (root) of any term can change due
to \union{} operations is bounded by $R_{\max}$.
\end{lemma}

\begin{proof}
  Each time a \union{} operation changes the parent of a root node, it is because two
  trees are being merged (non-root nodes are unaffected by \union{}). The rank of a root
  node increases only when two roots of equal rank are united, and as the maximum
  rank that can occur is $R_{\max}$, the \union{} operation can change the representative of
  any term at most $R_{\max}$ times.
\end{proof}

According to Lemma \ref{lemma:change-representative}, each term can
involved in at most $R_{\max}$ union operations that change its rank. Therefore,
all terms have stabilized with respect to their parent nodes after at most
$R_{\max} \in \Oh(\log n)$ iterations.


\subsection{Complexity Analysis}

Combining the complexities of the three phases of the algorithm, we obtain
the overall time complexity:
\[ \Oh(n) + \Oh(m \cdot \alpha(n)) + \Oh(n \cdot k \cdot \alpha(n) \log n) = \Oh(n \cdot k \cdot \alpha(n) \log n) \]

Observe that the term $k$ can be considered constant for a fixed signature $\Sigma$. As the
inverse Ackermann function $\alpha(n)$ grows extremely slowly with $n$ it can be considered
a constant for all practical purposes. This gives us the final result of a runtime of $\Oh(n \log n)$
for the entire algorithm.
